% Person 2 report section for Milestone 2. Include in the final M2 report once
% Person 1 supplies the exact Setup A parameters and run.

\section{Solr Query Configuration: Setup B}

Setup B uses eDisMax to weight fields and close phrase matches according to each
information need. Review-based needs query the nested \texttt{review} children and
use Solr's block-join parent parser with \texttt{score=max}; the returned documents
are parent movies ranked by their best matching review. Plot-based needs search the
parent movie fields directly. The output contract is one row per movie, with the
numeric TMDB \texttt{movie\_id} as the run and qrels identifier and Solr's
\texttt{id=movie-<movie\_id>} retained for traceability.

For Q1 and Q2, eDisMax searches \texttt{review\_text}; phrase boosts reward nearby
reviewer descriptors. The parent genre filter restricts Q1 to comedies and Q2 to
romance movies. For Q3, the overview receives the strongest boost because it states
the central plot, while the longer plot field provides supporting evidence; a
science-fiction genre boost raises likely matches without imposing a hard genre
filter. For Q4, the overview and plot fields are searched for FBI and serial-killer
language, and crime/thriller genres receive a secondary boost. For Q5, a required
cast-name clause identifies Robert De Niro while overview and plot text describe
the Mafia subject; cast and phrase boosts reward close matches.

The adopted Setup B run returned 186, 185, 60, 11, and 13 documents for Q1--Q5,
respectively. These are retrieval totals, not relevance counts. Manual inspection
found several false positives among the broadened Q3--Q5 candidates; pooled judging
will show whether the recall increase is useful. Full per-query results and parameters are saved in
\texttt{runs/B\_Q1.json} through \texttt{runs/B\_Q5.json}.

\begin{table}[t]
\centering
\caption{Planned baseline and improved retrieval configurations. Setup A details
must be confirmed against Person 1's implementation before final comparison.}
\label{tab:setup-ab}
\small
\begin{tabular}{p{1.15cm}p{2.4cm}p{4.5cm}}
\hline
Dimension & Setup A (baseline) & Setup B (improved) \\
\hline
Parser & Planned standard parser; verify exact parser with Person 1. & eDisMax for all five needs; Q1--Q2 use eDisMax inside the child query and Solr's parent block-join parser outside. \\
Fields & Planned single catch-all text field; verify indexed field with Person 1. & Q1--Q2: \texttt{review\_text}; Q3--Q4: overview (boost 3) and plot (boost 2); Q5: required \texttt{cast.name} clause plus cast, overview, and plot boosts. \\
Boosting & Default BM25 ranking is planned; confirm against the actual baseline. & Query-field boosts, phrase-field boosts, phrase slop, genre boosts, and query-specific minimum-match settings. \\
Returned unit & One movie per result. & One parent movie per result; review matches use maximum child score. \\
Run identifier & Numeric \texttt{movie\_id}. & Numeric \texttt{movie\_id}; Solr key \texttt{movie-<movie\_id>} is retained separately. \\
Depth & 20 results per query. & 20 results per query. \\
\hline
\end{tabular}
\end{table}

The frozen parameters are stored in \texttt{queries/setup\_b/Q1.json} through
\texttt{Q5.json}. Running \texttt{scripts/run\_setup\_b.py} writes one
\texttt{runs/B\_Qn.json} file per need and a TREC-format \texttt{runs/B.run}.
The JSON files retain the exact parameters, total number of matches, and the top
20 ranked movies. No Setup A run or relevance judgments are included here, so this
section makes no comparative performance claim and reports no evaluation metrics.

